\textbf{Performance-Counter IDs}\par
RDPMC selects a 64-bit counter with its 16-bit immediate ID. The three standard counters are mandatory and wrap
modulo $2^{64}$. Cold reset and warm RESET clear them to zero. PMC.EN controls whether they increment.

\Needspace{1.25in}
\manualtablecaption{Standard Performance Counters}
\begingroup\footnotesize
\setlength{\aboverulesep}{0pt}
\setlength{\belowrulesep}{0pt}
\setlength{\extrarowheight}{1.2pt}
\begin{center}
\begin{tabularx}{\linewidth}{@{}>{\raggedright\arraybackslash}p{0.55in}>{\raggedright\arraybackslash}p{1.05in}>{\raggedright\arraybackslash}X@{}}
\toprule
\rowcolor{ManualHeaderFill}
\textbf{ID} & \textbf{Name} & \textbf{Increment event}\\
\midrule
\texttt{1} & \texttt{CYCLE} & each implementation cycle while PMC.EN is set\\
\texttt{2} & \texttt{INSTRET} & each successfully committed architectural instruction\\
\texttt{3} & \texttt{PTWALK} & start of a hardware page walk after a paging-enabled TLB miss\\
\bottomrule
\end{tabularx}
\end{center}
\endgroup

Each successfully committed repeat-body iteration contributes one INSTRET event. A faulting instruction does not.
PTWALK includes a walk that later reports a page fault, but excludes a segment-bound or canonical-address failure
detected before a walk begins. ID zero is unavailable because index zero of the CPUID performance-counter leaf is its
common discovery header. Other IDs beyond the standard set are implementation-defined and are discovered through that
leaf.

\begin{manualraggedblock}
Every privilege level may read every implemented counter. Clearing PMC.EN freezes the current counter values
without clearing them.
\end{manualraggedblock}
